\section{Versão LaTeX convertida do documento-base}

\subsection{Escopo do projeto}

O projeto implementa uma plataforma de streaming ao vivo com backend em Java, frontend em TypeScript e orquestração por Docker Compose. O objetivo é comparar servidores RTMP e transcodificadores em uma arquitetura capaz de gerar e servir múltiplas variantes HLS, mantendo rastreabilidade entre execução e análise.

\subsection{Tecnologias principais}

\begin{table}[H]
\centering
\caption{Tecnologias que compõem a plataforma}
\begin{tabularx}{\textwidth}{@{}lXX@{}}
\toprule
\textbf{Camada} & \textbf{Tecnologia} & \textbf{Uso} \\
\midrule
Backend & Spring Boot / Spring Data JPA / Spring AMQP / Spring Retry & Controle da aplicação e integrações \\
Frontend & React / Vite / TanStack Query / HLS.js / Plyr & Interface de reprodução e observação \\
Ingestão RTMP & Nginx-RTMP ou SRS & Recebimento da live \\
Transcodificação & FFmpeg ou GStreamer & Geração de variantes HLS \\
Infraestrutura & Docker Compose & Isolamento e repetição do experimento \\
\bottomrule
\end{tabularx}
\end{table}

\subsection{Fluxo de streaming}

\begin{enumerate}[leftmargin=1.2cm]
  \item O streamer publica a live via RTMP.
  \item O servidor RTMP aciona callbacks no backend.
  \item O transcodificador gera a playlist mestre e os segmentos.
  \item Os espectadores consomem o fluxo via HLS.
  \item Os scripts registram os dados em CSV e logs versionados.
\end{enumerate}

\subsection{Uso dos diagramas}

Os diagramas originais permanecem mantidos em \texttt{docs/tcc/diagrams} como fonte Mermaid. Nesta primeira versão em LaTeX, eles podem ser exportados futuramente para PNG ou PDF e incluídos diretamente na compilação final. A estrutura textual, porém, já está preparada para receber as figuras sem alterar o restante do artigo.

\subsection{Matriz comparativa}

As quatro combinações estudadas são Nginx + FFmpeg, Nginx + GStreamer, SRS + FFmpeg e SRS + GStreamer. Esse recorte é suficiente para avaliar como a escolha do servidor e do transcodificador afeta o desempenho percebido em diferentes níveis de audiência.